% part: intuitionistic-logic
% chapter: propositions-as-types

\documentclass[../../../include/open-logic-chapter]{subfiles}

\begin{document}

\olchapter{int}{pty}{সূত্র-হিসেবে-টাইপ}

\begin{editorial}
  কারি--হাওয়ার্ড অনুরূপতা নিয়ে এই অধ্যায়টি
  \emph{অত্যন্ত পরীক্ষামূলক} খসড়া।
  এতে আরও ব্যাখ্যা ও প্রেরণা দরকার;
  সম্ভবত ভুল ও বাদ-পড়া অংশও আছে।
  স্বাভাবিকীকরণের প্রমাণটি পুনর্বিবেচনা
  ও বিস্তৃত করা উচিত। গুণন-টাইপের
  কোনো উদাহরণ নেই। বিন্যাস ও
  সরলীকরণ রূপান্তর এখানে আলোচিত নয়।
  (টাইপযুক্ত) ল্যাম্বডা ক্যালকুলাস নিয়ে
  উপাদান যোগ হলে অধ্যায়টির অর্থ
  অনেক বেশি স্পষ্ট হবে; এখানে মূলত
  সেই জ্ঞানই পূর্বধরা হয়েছে।
  অত্যন্ত সতর্কতার সঙ্গে ব্যবহার করুন।
\end{editorial}

\olimport{introduction}
%\olimport{sequent-natural-deduction}
\olimport{proof-terms}
\olimport{proofs-to-terms}
\olimport{terms-to-proofs}
\olimport{types}
\olimport{reduction}
\olimport{type-preservation}
\olimport{normalization}

\OLEndChapterHook

\end{document}
